%!TEX root = ../thesis.tex
%*******************************************************************************
%****************************** Chương 6 ***************************************
%*******************************************************************************

\chapter{KẾT LUẬN}
\label{chap:conclusion}

\ifpdf
    \graphicspath{{Chapter6/Figs/Raster/}{Chapter6/Figs/PDF/}{Chapter6/Figs/}}
\else
    \graphicspath{{Chapter6/Figs/Vector/}{Chapter6/Figs/}}
\fi

%********************************** Mục 6.1 ************************************
\section{Tóm tắt}
\label{sec:con-summary}

Khóa luận nghiên cứu cách thích nghi một vision-language model mã nguồn mở cho
hỏi đáp trên ảnh tiếng Việt dưới ràng buộc tài nguyên chặt: chỉ dùng
parameter-efficient fine-tuning và chỉ một GPU Colab. Nhóm chọn
Qwen2.5-VL-3B-Instruct làm baseline, tái lập nó trên TextVQA val (79,60 so với
79,3 công bố), đo baseline zero-shot trên ViVQA và ViTextVQA, rồi thích nghi bằng
QLoRA với 7,37 triệu tham số huấn luyện, tức 0,196\% tổng số tham số.

Cải tiến đề xuất là QLoRA nhận biết OCR: chữ trong ảnh được đọc bằng VietOCR,
sắp theo thứ tự đọc và đưa vào prompt ngay trong lúc huấn luyện. Trên toàn bộ
10.028 câu test ViTextVQA, phương pháp này đạt EM 53,35, hơn QLoRA thông thường
với cùng dữ liệu 3,09 điểm ($p = 6{,}7\times10^{-21}$) và cao hơn kết quả công bố
tốt nhất trước đây (EM 25,48). Trên ViVQA, QLoRA nâng EM từ 22,09 lên 76,64.

%********************************** Mục 6.2 ************************************
\section{Trả lời các câu hỏi nghiên cứu}
\label{sec:con-rq}

\begin{description}
  \item[RQ1.] \textit{Zero-shot đạt bao nhiêu, và lỗi tiếng Việt nào chiếm ưu thế?}
  Zero-shot đạt EM 22,09 trên ViVQA và 29,12 trên ViTextVQA. Lỗi chiếm ưu thế là
  sai định dạng: mô hình trả lời dài hơn đáp án (2,9 so với 1,6 từ trên ViVQA) và
  đôi khi trả lời bằng tiếng Anh. Đáp án có dấu thanh khó hơn rõ rệt so với đáp
  án chỉ gồm chữ Latin. Prompt few-shot định dạng không sửa được các lỗi này
  (H3 không được ủng hộ).

  \item[RQ2.] \textit{QLoRA cải thiện bao nhiêu khi không dùng full fine-tuning?}
  QLoRA tăng EM thêm 21,1 điểm trên ViTextVQA và 54,6 điểm trên ViVQA (mô hình
  chính thức C3), chỉ với 0,196\% tham số, một GPU L4, tối đa 12,75 giờ huấn luyện
  và VRAM đỉnh dưới 8,4~GB. Mức tăng chủ yếu đến từ việc học đúng định dạng đáp án
  và không do rò rỉ ảnh: trên tập test sạch, mức tăng vẫn là 44,1 điểm (H1 được
  ủng hộ).

  \item[RQ3.] \textit{Chữ OCR có giúp, và có cộng hưởng với QLoRA không?}
  Có, nhưng chỉ khi mô hình được huấn luyện cùng OCR. Thêm OCR vào prompt của một
  mô hình chưa từng thấy OCR không cải thiện; huấn luyện cùng OCR cho hiệu ứng tương
  tác $+4{,}9$ điểm EM, ổn định trên ba seed, và lợi ích không giảm khi tăng dữ
  liệu. Mức tăng ANLS ($+1{,}87$) thấp hơn ngưỡng 3 điểm đặt trước, nên H2 được
  ủng hộ một phần: OCR giúp trả lời đúng nguyên chuỗi nhiều hơn là giúp gần đúng
  theo ký tự.

  \item[RQ4.] \textit{Sau khi thích nghi, mô hình bịa nội dung nhiều hơn hay ít hơn?}
  Trên POPE-vi, B3-full không tăng hallucination theo luật đặt trước ($\Delta$F1
  từ $+2{,}5$ đến $+5{,}5$), nhưng tỉ lệ báo sai vật thể tăng ở cách chọn
  adversarial. C3 không đo được vì đã quên cách trả lời có/không. Duyệt tay cho
  thấy hallucination chỉ chiếm khoảng 2,3\% (B3-full) và 3,7\% (C3) mọi câu trả
  lời. H4 được ủng hộ một phần.
\end{description}

%********************************** Mục 6.3 ************************************
\section{Đóng góp}
\label{sec:con-contrib}

\begin{enumerate}
  \item \textbf{QLoRA nhận biết OCR cho VQA tiếng Việt}, với cơ chế được tách bạch
  bằng thí nghiệm $2\times2$: lợi ích đến từ việc dạy mô hình dùng OCR, không đến
  từ việc thêm OCR vào prompt.
  \item \textbf{Thích nghi hiệu quả và tái lập được.} Kết quả trên ViTextVQA cao
  hơn kết quả công bố trước đây với 0,196\% tham số trên một GPU; mọi con số có
  log, seed và cấu hình, kèm gói tái lập gồm mã nguồn, adapter, cấu hình, script
  suy luận và hướng dẫn chạy lại.
  \item \textbf{Đo các vấn đề riêng của tiếng Việt.} Bốn vấn đề V1--V4 đều có số
  đo cụ thể; POPE-vi là bộ đo hallucination tiếng Việt gồm 4.985 câu trên 500 ảnh
  chưa thấy khi huấn luyện.
  \item \textbf{Phân tích lỗi hai tầng.} Phân tích tự động theo kiểu lỗi và trạng
  thái OCR, đối chiếu với duyệt tay mù, cho thấy EM đánh giá thấp mô hình và ước
  lượng phần còn có thể cải thiện nhờ OCR tốt hơn (khoảng 17 điểm EM).
  \item \textbf{Prototype đã kiểm chứng.} Ứng dụng web trả lời kèm chữ OCR và mức
  tin cậy, được kiểm cho cùng câu trả lời với đường đánh giá chính thức, với độ
  trễ trung vị khoảng 1 giây trên một GPU L4.
\end{enumerate}

%********************************** Mục 6.4 ************************************
\section{Giới hạn}
\label{sec:con-limits}

Các giới hạn chính đã được phân tích ở mục~\ref{sec:dis-validity}: mô hình chính
chỉ chạy một seed; EM quá khắt khe với câu trả lời đúng ý nhưng khác cách viết;
chất lượng OCR tiếng Việt cho chữ cảnh là nút thắt; ngưỡng ANLS của H2 chưa đạt;
F1 trên POPE-vi che mất tỉ lệ báo sai; so sánh với kết quả công bố phụ thuộc cách
chuẩn hóa đáp án của bài báo; và lần duyệt tay cuối chỉ có một người gán.

%********************************** Mục 6.5 ************************************
\section{Hướng phát triển}
\label{sec:con-future}

\begin{enumerate}
  \item \textbf{Nâng cấp OCR tiếng Việt cho chữ cảnh}: fine-tune VietOCR trên chữ
  cảnh và cải tiến thứ tự đọc. Đây là hướng có trần cải thiện lớn nhất theo phân
  tích lỗi.
  \item \textbf{OCR có vị trí}: đưa tọa độ và thứ tự các hộp chữ vào mô hình theo
  hướng của M4C~\cite{hu2020m4c}, để giảm lỗi cắt sai ranh giới cụm chữ.
  \item \textbf{Đánh giá có người chấm và thước đo chấp nhận đồng nghĩa}, vì hơn
  một nửa lỗi EM là đúng ý.
  \item \textbf{Cân bằng câu hỏi có/không khi huấn luyện} để tránh quên định dạng
  như C3, và đo hallucination chữ trên ảnh có chữ, không chỉ hallucination vật thể.
  \item \textbf{Hiệu chỉnh độ tin cậy} (ví dụ temperature scaling) để prototype có
  thể hiện xác suất thay cho mức.
  \item \textbf{Thêm seed cho mô hình chính} và thử các mô hình nền mới hơn dưới
  cùng giao thức đánh giá.
\end{enumerate}

Kết quả của khóa luận cho thấy một nhóm có tài nguyên hạn chế vẫn có thể thích
nghi một VLM lớn cho tiếng Việt và đạt kết quả vượt các công trình trước, miễn là
mỗi thay đổi được đo dưới cùng điều kiện và mỗi kết luận được kiểm định.
